\vspace{-1pt}

\section{Related Work}
\vspace{-1pt}


\noindent \textbf{Vision-Language Models for Long Sequence Understanding.}\;\; 
Early Vision-Language Models (VLMs), such as GPT-4V and Gemini-1.5~\cite{openai2024gpt4ocard, geminiteam2024gemini15unlockingmultimodal}, showcased powerful multimodal reasoning by integrating visual encoders with large language models. Open-source efforts like Llama-Vid~\cite{li2024llama}, IDEFICS~\cite{idefics2023}, VideoChat~\cite{li2023videochat}, Video-LLaMA~\cite{cheng2024videollama}, and others~\cite{li2024videochat, lin2023video, wang2022internvideo, wang2024internvideo2, li2024llava, luo2023valley, ataallah2024minigpt4} have further advanced capabilities, often matching or exceeding proprietary systems on various benchmarks. While effective on static images or short video clips, scaling VLMs to long videos (minutes to hours) introduces a significant challenge. The core difficulty lies in managing the immense volume of visual data, which generates an excessive number of tokens that quickly exceed the typical context windows and computational capacities of the LLM components during both training and inference. Current research addressing long-context VLMs primarily explores two directions: (i) methods that significantly extend the effective context length of the underlying transformer architectures, such as LongVA, LongVILA, LongViTA, and LLaVA-NeXT~\cite{zhang2024long, chen2024longvila, shen2025long, zhang2024llavanext}, and (ii) strategies focused on reducing the number of visual tokens fed to the LLM via selection or compression modules~\cite{korbar2024text, jin2024chat, shu2024video, song2024moviechat}. While context extension allows processing more tokens, it often incurs high computational costs. Conversely, existing visual compression techniques, while reducing token count, frequently rely on heuristics or fixed operations, or uniform sampling across the video~\cite{li2023videochat, cheng2024videollama, zhang2024long, li2024llava, li2024videochat, ataallah2024minigpt4} that struggle to preserve critical temporal and fine-grained information, particularly when pushed to high compression ratios necessary for very long videos. Other orthogonal approaches tackle long videos by processing shorter segments independently and employing retrieval mechanisms~\cite{ataallah2024goldfish}.

\noindent \textbf{Video Token Compression in Multimodal LLMs.}\;\; 
Given the computational constraints of processing dense video streams with VLMs, compressing visual tokens is a vital approach. Various methods have been proposed to distill frame sequences into more manageable representations. These include temporal pooling strategies~\cite{fei2024video, maaz2023video}, heuristic or saliency-based frame selection often guided by simple visual priors or attention mechanisms~\cite{shu2024video, song2024moviechat}, and lightweight learned modules or bottlenecks aimed at reducing the token sequence, such as Llama-Vid~\cite{li2024llama} which explores compressing frames potentially to a single token by leveraging cross-attention with a text query. Relatedly, techniques for dynamic token pruning or merging have been explored for both images and video~\cite{ma2023image, xu2022groupvit, bolya2022token, lee2024multi, ren2023testa, choi2024vid, jin2024chat, xu2024slowfast}. However, a critical limitation of many existing approaches, especially those relying purely on heuristics or fixed pooling, is that they are not trained to consolidate the information from dropped tokens into the retained ones within the LLM's representation space. Aggressively applying these methods to achieve the extreme compression ratios necessary for very long videos risks discarding vital visual details, thus hindering deep understanding. Furthermore, while frame selection is intuitive, applying it effectively across extremely long sequences without being hampered by issues like internal LLM position biases remains challenging. Therefore, developing methods that can perform learned and highly efficient visual compression, potentially towards representing each frame with an unprecedentedly low token count while preserving critical information, remains a significant research frontier.